\subsection{All Quadratic Functions are Power Functions} 
We saw before that all graphs of quadratic functions have the same shape, a \makeblank[1.5in]. This suggests a family of functions!
\begin{fact}[Claim: All Quadratic Functions are Power Functions]
All quadratic functions can be written in the form
$$f(x)=a(x-h)^2+k$$
for real numbers $a$, $h$, and $k$. This form is called the \term{vertex form} of the parabola.
\end{fact}
We will prove this claim by developing a procedure that will let us write any quadratic function in this form. However, note that this isn't true of all polynomials.
\begin{Example}
The graphs of three polynomial functions of degree 4 are shown below. (Note the unlabeled axes; we are interested in the general shape of the graph.)
\begin{center}\begin{tikzpicture}[x=0.5\linespace,y=0.5\linespace]
\setgraphsquare{5}
\begin{scope}[xshift=-0.3\linewidth]
\AxesBlank
\draw[graph,domain={-7^(1/4)}:{7^(1/4)}] plot (2*\x,{pow(\x,4)-2});
\draw (0,-5.5) node[yshift=-2ex,anchor=base,font=\small] {$f(x)=x^4-2$};
\end{scope}
\begin{scope}[xshift=0]
\AxesBlank
\draw[graph,domain=-3.53:3.02] plot (\x,{0.1*pow(\x-1,3)*(\x+3)});
\draw (0,-5.5) node[yshift=-2ex,anchor=base,font=\small] {$g(x)=x^4-6x^2+8x-3$};
\end{scope}
\begin{scope}[xshift=0.3\linewidth]
\AxesBlank
\draw[graph,domain={-sqrt((9+sqrt(181))/2)}:{sqrt((9+sqrt(181))/2)},samples=50] plot (\x,{0.2*(pow(\x,4)-9*pow(\x,2))});
\draw (0,-5.5) node[yshift=-2ex,anchor=base,font=\small] {$h(x)=x^4-9x^2$};
\end{scope}
\end{tikzpicture}\end{center}
We can see that these have very different \makeblank[1in], so these functions to \emph{not} form a family!
\end{Example}